\input{preamble.tex}

% --------------------------------------------------------------------
% --------------------------------------------------------------------
% ---- ENTER YOUR PROJECT DETAILS BELOW

% ---- title of the project
\title{SHORT TITLE OF YOUR PROJECT}

% ---- names of team members
\author{
    \small Name1 Surname1 \and
    \small Name2 Surname2 \and
    \small Name3 Surname3 \and
    \small Name4 Surname4
    }

% ---- enter the team number below
\rhead{Team Number XX}

% --------------------------------------------------------------------
% --------------------------------------------------------------------


% --------------------------------------------------------------------
% --------------------------------------------------------------------
% ---- DO NOT CHANGE THIS BLOCK OF CODE
\begin{document}
\maketitle
% --------------------------------------------------------------------
% --------------------------------------------------------------------



% --------------------------------------------------------------------
% --------------------------------------------------------------------
% ---- WRITE YOUR REPORT BELOW

\section{Introduction}
\lipsum[1]

Here is an example of an in-text
citation~\citep{silva-rodriguez_conformal_2025}. And this is an example where
you want to use the author's name in the text, where \citet{fillioux_are_2025}
have done something interesting.


\section{Related Work}
\lipsum[3]

\section{Our contribution}
\lipsum[4][1-15]

\section{Methods}
\subsection{Data}
\lipsum[1][1-5]

\subsection{Name of main method/theory/idea}
\lipsum[2][1-10]

You can use the \verb|align| environment to manage equations.
\begin{align}
    \mathbf{A} \mathbf{X} & =  \lambda \mathbf{X},  \textrm{and} \label{eq:e1} \\
    \dot{\mathbf{X}}  & =  \frac{\mathrm{d}\mathbf{X}}{\mathrm{dt}} \label{eq:e2}  
\end{align}
are widely used in deriving models from data. And then you can use
cross-referencing to refer to equations in the text, like Eq.~\ref{eq:e1}, and
Eq.~\ref{eq:e2}, provided you have used \verb|\label| in each equation.

\section{Results}
\subsection{First results: Plots}

This is a cross-reference, as seen in Fig.~\ref{fig:main_result} to the figure
in the text. Figures have to be placed either at the top of the page or at the
bottom of the page, and if in case they do not fit in any of these
arrangements, they have to be placed in a page of their own.

\lipsum[5]

\begin{figure}[tbp]
    \begin{center}
        \includegraphics[width=0.75\textwidth]{./figs/result.png}
        \caption{This is a first results figure.}
        \label{fig:main_result}
    \end{center}
\end{figure}

\subsection{Second results: Dashboard}
\lipsum[4]

\subsection{Third results: Use case}
\lipsum[3]

\section{Conclusion and Outlook}
\lipsum[7][1-20]

\section{Acknowledgements}
\lipsum[8][1-5]

\section{Author Contributions}
\lipsum[1][1-9]

\section{Code availability}
The code for this project is available at the following GitHub repository:
\url{https://github.com/bedartha/cl-todo}.

% --------------------------------------------------------------------
% --------------------------------------------------------------------
% ---- DO NOT CHANGE THIS BLOCK OF CODE
\section{References}
\bibliography{refs}
% --------------------------------------------------------------------
% --------------------------------------------------------------------


\section{Appendix}

\subsection{Additional figures}
If you have more figures which do not form part of the main narrative, you can
dump them here.

\subsection{Algorithms}

You can dump algorithms here as well.

\begin{minipage}{0.7\textwidth}
    \begin{algorithm}[H]
        \caption{An algorithm with caption}\label{alg:cap}
        \begin{algorithmic}
        \Require $n \geq 0$
        \Ensure $y = x^n$
            \State $y \gets 1$
            \State $X \gets x$
            \State $N \gets n$
            \While{$N \neq 0$}
                \If{$N$ is even}
                    \State $X \gets X \times X$
                    \State $N \gets \frac{N}{2}$  \Comment{This is a comment}
                \ElsIf{$N$ is odd}
                    \State $y \gets y \times X$
                    \State $N \gets N - 1$
                \EndIf
            \EndWhile
        \end{algorithmic}
    \end{algorithm}
\end{minipage}


\subsection{Code snippets}

Finally, you can also showcase a few important code snippets which you think
are really important so as not to wait till the reader actually opens the
Github repo. But to run this, you will need to compile the \verb|tex| document with
the \verb|-shell-escape| flag. This means that instead of running:\\
\verb|$ pdflatex report.tex|\\
as you would usually do, you have to now run\\
\verb|$ pdflatex --shell-escape report.tex|\\
instead.

\begin{minted}[bgcolor=white, fontsize=\small, linenos, frame=none]{python}
# Function to check prime numbers
def is_prime(n):
    if n < 2:
        return False
    for i in range(2, int(n**0.5) + 1):
        if n % i == 0:
            return False
    return True

# Print primes from 1 to 20
for i in range(1, 21):
    if is_prime(i):
        print(f"{i} is prime")
\end{minted}


% --------------------------------------------------------------------
% --------------------------------------------------------------------
% ---- DO NOT CHANGE THIS BLOCK OF CODE
\end{document}
% --------------------------------------------------------------------
% --------------------------------------------------------------------
